\documentclass[11pt]{article}

% --------------------------------------------------
% Packages
% --------------------------------------------------
\usepackage[margin=1in]{geometry}
\usepackage{amsmath, amssymb, amsthm}
\usepackage{mathtools}
\usepackage{graphicx}
\usepackage{float}
\usepackage{enumitem}
\usepackage{hyperref}
\usepackage{xcolor}
\usepackage{tikz}
\usetikzlibrary{arrows.meta, positioning, calc}

% Pseudocode and algorithms
\usepackage{algorithm}
\usepackage{algpseudocode}

% Better fonts (optional)
\usepackage{lmodern}

% --------------------------------------------------
% Hyperref setup
% --------------------------------------------------
\hypersetup{
    colorlinks=true,
    linkcolor=blue,
    citecolor=blue,
    urlcolor=blue
}

% --------------------------------------------------
% Theorem-like environments
% --------------------------------------------------
\theoremstyle{definition}
\newtheorem{definition}{Definition}[section]
\newtheorem{theorem}{Theorem}[section]
\newtheorem{claim}{Claim}[section]
\newtheorem{remark}{Remark}[section]
\newtheorem{recall}{Recall}[section]

\theoremstyle{proof}
% Proof environment is built-in via amsthm

% --------------------------------------------------
% Custom commands
% --------------------------------------------------
\newcommand{\Gen}{\mathsf{Gen}}
\newcommand{\Enc}{\mathsf{Enc}}
\newcommand{\Dec}{\mathsf{Dec}}
\newcommand{\Sig}{\mathsf{Sig}}
\newcommand{\Ver}{\mathsf{Ver}}
\newcommand{\negl}{\mathsf{negl}}

% --------------------------------------------------
% Document
% --------------------------------------------------
\begin{document}

\title{Public Key Crypto notes - 11}
\author{By contributors}
\date{April 27. 2026}
\maketitle

\section{Projective coordinates}

\begin{recall}
	$\mathbb{R}^2$ is the two-dimensional affine plane. In this space, there are:
	\begin{itemize}
		\item Points,
		\item Lines
		\item $\mathbb{P}(\mathbb{R}^2)$ is the two-dimensional projective space, that contains:
		\begin{itemize}
			\item affine points,
			\item affine lines,
			\item a "projective" point: meeting points of parallel lines,
			\item a "projective" line: line of all projective points.
		\end{itemize}
	\end{itemize}
\end{recall}

\begin{definition}[Projective points]
	For $x1, y_1, z_1, x_2, y_2, z_2$, where $((x_i, y_i, z_i) \ne (0, 0, 0))$, we say $(x_1, y_1, z_1) \sim (x_2, y_2, z_2)$ if $x_2 = \lambda x_1, y_2 = \lambda y_1, z_2 = \lambda z_1$ for some $\lambda \ne 0$.
	
	This $\sim$ is an equivalence relation. Let $(x:y:z)$ be a representative of the equivalence classes of $(x, y, z)$.
	
	Example: $(1:1:1) = (2:2:2) = (3:3:3)$, but $(1:1:1) \ne (2:1:0)$.
\end{definition}

\begin{remark}
	If $z \ne 0$, then $(x:y:z) = (\frac{x}{z}:\frac{y}{z}:1)$. These will be the affine points.
	
	If $z = 0$, then $(x:y:z) = (x:y:0)$. These will be the projective points.
\end{remark}

\begin{definition}[Homogenous polynomial]
	A polynomial $F \in \mathbb{R}[x, y, z]$ of form $F = \Sigma_{c_{i,j,k}}x^i y^j z^k$ is \emph{homogenous} if $c_{i, 0, k} \ne 0$, then $i + j + k = d$ for some $d$. Here, $\deg F = d$.
	
	Examples: $F_1 = c_1 x + c_2 y + c_3$ is not homogenous, while $F_2 = c_1 x + c_2 y + c_3 z$ is $\deg = 1$ homogenous.
\end{definition}

\begin{remark}
	If $F$ is homogenous and $(x_1:y_1:z_1) = (x_2:y_2:z_2)$, then $F(x_1,y_1,z_1) = 0 \iff F(x_2, y_2, z_2) = 0$.
	
	Indeed, if:
	
	\begin{align}
		x_2 &= \lambda x_1, \qquad
		y_2 = \lambda y_1, \qquad
		z_2 = \lambda z_1. \\
		\intertext{Then}
		F(x_2,y_2,z_2)
		&= F(\lambda x_1,\lambda y_1,\lambda z_1) \\
		&= \sum_{i+j+k=d} c_{i,j,k} (\lambda x_1)^i (\lambda y_1)^j (\lambda z_1)^k \\
		&= \lambda^d \sum_{i+j+k=d} c_{i,j,k} x_1^i y_1^j z_1^k \\
		&= \lambda^d F(x_1,y_1,z_1).
	\end{align}
\end{remark}

\begin{definition}
	A line of the projective space is the zero-set of a $\deg = 1$ homogenic polynomial.
	
	Example 1: Take a horizontal line in $\mathbb{R}^2$. $y + c = 0 \rightarrow$ we can homogenize this: $y + cz = 0$.
	
	$(x:y:z)$ is a zero.
	
	\begin{itemize}
		\item If $z = 1 \rightarrow$ affine point.
		\item If $z = 0 \rightarrow y = 0 \rightarrow x \ne 0$.
	\end{itemize}
	
	So $(x:0:0) = (1:0:0)$.
	
	Moreover, $(1:0:0) = (-1:0:0)$.
	
	Example 2: Take a vertical line $x + c = 0 \rightarrow$ make it homogenous $\rightarrow x + cz = 0 \rightarrow x + 0y + cz = 0$.
	
	\begin{itemize}
		\item If $z = 1 \rightarrow$ affine points.
		\item If $z = 0 \rightarrow x = 0 \rightarrow y \ne 0 \rightarrow (0:y:0) = (0:1:0)$.
	\end{itemize}
\end{definition}

\subsection{Nice subsection title}

Let $E:y^2 = x^3 + Ax + B$ and $F = x^3 + Ax + B - y^2$. Let's homogenize $F$:

$H = x^3 + Axz^2 + Bz^3 - y^2z$

The point of E:

\begin{itemize}
	\item If $z = 1 \rightarrow H(x, y, 1) = 0 \iff (x, y) \in E$.
	\item If $z = 0 \rightarrow H(x, y, 0) = x^3 \rightarrow x = 0 \rightarrow y \ne 0$.
\end{itemize}

Thus, $(x:y:0) = (0:1:0)$ is the only fully projective point of $E$.

\begin{remark}\leavevmode
	\begin{itemize}
		\item Similar addition low can be given, but we do not have to distinguish $P = \infty, P \ne \infty$.
		\item In implementations, only projective coordinates are used.
		\item It speeds up computations as modular inversions can be eliminated.
	\end{itemize}
\end{remark}

\section{Elliptic curve cryptography}

\begin{remark}
	Let $P$ be a prime, and let $E$ be an elliptic curve over $\mathbb{Z}_p$. Then, the computation $n \rightarrow nP$ can be easily computed using the double-and-add algorithm (where $n$ is an integer, and $P \in E$).
	
	i.e. $n = \sum_{i=0}^{\lambda} b_i z^i, b_i \in \{0, 1\}$.
	
	$nP = 2(2((\dots (2b_iP)\dots) \oplus b_2P) \oplus b_1P) \oplus b_0P$.
	
	There are $\approx 2l = 2\log_2n$ many curve operations.
	
	But computing $n$ from $P, nP$ is hard. This is called the DLog problem on the curve.\footnote{Where previously, we had $g^n$ as the basis of security, we now have $nP$ for the same purpose.}
\end{remark}

Techniques for attack:

\begin{enumerate}
	\item Brute-force: trying all possible $n$ values. The possible values of $n \approx |E| \approx P$ (if $P$ is not trivial).
	\item There are smarter algorithms (like baby-step-giant-step) having a running time of approximately $\sqrt{P}$.
	\item Index calculus does \emph{not} work.
\end{enumerate}

Hence, to achieve the same security level AES-128, one needs $P$ to be approximately $2^{256}$ bit size. For RSA and ElGamal encryption, one needs $3000$ bits for the same security level.

\subsection{Elliptic Curve Diffie-Hellman key exchange}

Let $P$ be a prime and $E$ be an elliptic curve over $\mathbb{Z}_P, P \in E$ of order $l$ (i.e. $lP = \infty$, for $0 < k < l, kP + \infty$).

As $E \stackrel{\sim}{=} \mathbb{Z}_N$ or $\mathbb{Z}_N \times \mathbb{Z}_M$, $E$ and $P$ can be chose such that $l \approx P$.

\begin{enumerate}
	\item Alice chooses $a \in_R \mathbb{Z}_l$ and sends it to Bob.
	\item Bob chooses $b \in_R \mathbb{Z}_l$ and sends it to Alice.
	\item $a(bP) = abP = b(aP)$.
	\item The shared key is $x(abP)$
\end{enumerate}

\subsection{Elliptic Curve ElGamal encryption}

\begin{itemize}
	\item $\Gen$:
	\begin{itemize}
		\item Let $p$ be a prime and $E$ be an elliptic curve over $\mathbb{Z}_P, P \in E$ or order $l$.
		\item $a \in_R \mathbb{Z}_l, Q = aP$.
		\item $sk = a, pk = (p, E, P, Q)$.
	\end{itemize}
	\item $\Enc$:
	\begin{itemize}
		\item Given $pk = (p, E, P, Q)$ and let $M \in E$ be a message.
		\item $c = (rP, M \oplus rQ)$ with $r \in_R \mathbb{Z}_l$.
	\end{itemize}
	\item $\Dec$:
	\begin{itemize}
		\item Given $sk = a$, and the ciphertext $c = (R, S)$, the message is $M = S \ominus aR$.
	\end{itemize}
\end{itemize}

\end{document}